(1) 对速率 $v=\beta c$，洛伦兹因子 $\gamma=(1-\beta^2)^{-1/2}$。粒子固有时间增量与实验室时间满足 $d\tau=dt/\gamma$，故实验室观察到的寿命均值为 $\langle t\rangle=\gamma\tau_0$。这是时间膨胀而非衰变过程在粒子自身参照系内变慢。
(2) 固有时间 $\tau$ 内存活概率 $P(\tau)=e^{-\tau/\tau_0}$。穿越距离 $H$ 所需实验室时间为 $t=H/v$，对应固有时间 $\tau=H/(\gamma v)$，所以存活概率为 $P_{\rm surv}(H)=\exp[-H/(\gamma v\tau_0)]$，无需假定速度接近光速才可使用。
(3) 令 $P_{\rm surv}(H_{1/2})=1/2$，得 $H_{1/2}=\gamma v\tau_0\ln2=\gamma\beta c\tau_0\ln2$。平均衰变长度定义为 $\ell_{\rm decay}=\gamma v\tau_0$，对应存活概率降至 $1/e$ 的长度，因此半数存活距离为 $\ell_{\rm decay}\ln2$，两者不能混淆。
(4) 若错误忽略时间膨胀，得到 $P_{\rm naive}=\exp[-H/(v\tau_0)]$，高速时严重低估存活概率。当 $\beta\to0$ 时对固定 $H>0$ 两种概率均趋零；当 $\beta\to1$ 时 $\gamma\to\infty$，固定 $H$ 的存活概率趋一。真实穿过物质时粒子可能散射、损能量或被吸收，使 $v,\gamma$ 随位置改变，需用沿路径积分 $\exp[-\int ds/(\gamma(s)v(s)\tau_0)]$ 并另计非衰变损失。
